\begin{pspicture}(-2.4,-2.4)(5.2,2.4)
\psset{arrowsize=4pt 2}

% The four sites sit on a ring; the dashed circle is the chain of bonds and
% the wrap-around bond closes it.
\pscircle[linewidth=0.5pt,linestyle=dashed](0,0){1.5}
\psdot[dotsize=4pt](1.5,0)
\psdot[dotsize=4pt](0,1.5)
\psdot[dotsize=4pt](-1.5,0)
\psdot[dotsize=4pt](0,-1.5)

% One spin per site, drawn along z.  The alternating up/down pattern is the
% ordered (Neel) configuration the antiferromagnetic coupling favors.
\psline[linewidth=1pt]{->}(1.5,-0.4)(1.5,0.4)
\psline[linewidth=1pt]{->}(0,1.9)(0,1.1)
\psline[linewidth=1pt]{->}(-1.5,-0.4)(-1.5,0.4)
\psline[linewidth=1pt]{->}(0,-1.1)(0,-1.9)

\rput[r](1.0,0){\small $i=0$}
\rput[t](0,1.0){\small $i=1$}
\rput[l](-1.0,0){\small $i=2$}
\rput[b](0,-1.0){\small $i=3$}

\rput[bl](1.2,1.2){$J$}
\rput[br](-1.2,1.2){$J$}
\rput[tr](-1.2,-1.2){$J$}
\rput[tl](1.2,-1.2){$J$}

% The transverse field points along x, at a right angle to the spin axis.
\psline[linewidth=1pt]{->}(2.7,0)(4.1,0)
\rput[b](3.4,0.2){$h_x$}

% The axes: spins are measured along z, the field is applied along x.
\psline[linewidth=0.5pt]{->}(2.9,-2.0)(2.9,-1.2)
\psline[linewidth=0.5pt]{->}(2.9,-2.0)(3.7,-2.0)
\rput[b](2.9,-1.1){\small $z$}
\rput[l](3.8,-2.0){\small $x$}

\end{pspicture}
